\documentclass[../../main.tex]{subfiles}
\begin{document}

\begin{center}
  {\fontsize{48pt}{58pt}\selectfont \textbf{1998}\normalsize\textbf{年度}}
\end{center}

\begin{table}[htb]
  \centering
  \caption{出題テーマと難易度}
  \begin{tabular}{|c|p{9cm}|c|}
    \hline
    \textbf{問題番号} & \textbf{テーマ} & \textbf{難易度} \\
    \hline
    第1問 & 三次関数の極大値と極小値の差 & \\
    \hline
    第2問 & 空間の整数点の個数と極限 & \\
    \hline
    第3問 & 接する円の列と有理数の極限 & \\
    \hline
    第4問 & 床関数の値の種類と数え上げ & \\
    \hline
    第5問 & ベクトルの射影で定まる数列の収束 & \\
    \hline
    第6問 & 四角錐と円柱の外側部分の体積 & \\
    \hline
  \end{tabular}
\end{table}

\begin{table}[htb]
  \centering
  \caption{解答}
  \renewcommand{\arraystretch}{1.4}
  \begin{tabular}{|c|c|p{9cm}|}
    \hline
    \textbf{問題番号} & \textbf{小問} & \textbf{答え} \\
    \hline
    第1問 & -- & $a=\pm1$ \\
    \hline
    第2問 & -- & $\displaystyle\lim_{n\to\infty}
    \dfrac{f(n)}{n^3}=\dfrac{8}{3}$ \\
    \hline
    第3問 & (1) & $q_n$は整数（証明問題） \\
    \cline{2-3}
    & (2) & $p_n$も整数で，$p_n$と$q_n$は互いに素
    （証明問題） \\
    \cline{2-3}
    & (3) & 指定された不等式が成り立ち，
    $\displaystyle\lim_{n\to\infty}x_n
    =\dfrac{\sqrt{5}-1}{2}$ \\
    \hline
    第4問 & -- & $26n+1$個 \\
    \hline
    第5問 & -- & $\theta=0,\pi$，
    $\dfrac{\pi}{2}\le\theta\le\dfrac{5\pi}{6}$，
    $\dfrac{3\pi}{2}\le\theta\le\dfrac{11\pi}{6}$で収束．
    両区間の内部と$0,\pi$では極限$(0,0)$．
    $\theta=\dfrac{\pi}{2},\dfrac{3\pi}{2}$では
    $\left(1,-\sqrt{3}\right)$，
    $\theta=\dfrac{5\pi}{6},\dfrac{11\pi}{6}$では
    $\left(-\dfrac{1}{2},\dfrac{\sqrt{3}}{2}\right)$ \\
    \hline
    第6問 & -- & $V=4(\sqrt{2}+1)-3\pi$ \\
    \hline
  \end{tabular}
\end{table}

\paragraph{試験形態}
6問．試験時間は要確認．

\end{document}
